% ================================================================
% DISCUSSION
% ================================================================
%
% Source notes: PASSTHROUGH.md, IMPORT_LEAKAGE.md, B2B_LEAKAGE.md.
% Cross-reference §\ref{sec:passthrough}, §\ref{sec:import_leakage},
% §\ref{sec:b2b_leakage}.
%
% ----------------------------------------------------------------

% ----------------------------------------------------------------
% 6.1  PUTTING THE THREE FINDINGS TOGETHER
% ----------------------------------------------------------------
\subsection{Pass-Through Without Leakage}\label{sec:discussion:joint}

% TODO: develop the headline story.
%
% The three findings are mutually consistent and tell one story:
%
%   (a) Carbon costs are paid forward — Belgian regulated-NACE producer
%       prices rise with realized ETS exposure (§\ref{sec:passthrough}).
%
%   (b) Belgian buyers do not substitute toward unregulated foreign sources
%       (§\ref{sec:import_leakage}, contra CdGM France).
%
%   (c) Belgian buyers do not substitute toward unregulated domestic
%       suppliers either (§\ref{sec:b2b_leakage}).
%
% Implication: ETS pass-through to final / downstream buyers is real, but
% the supply-substitution channels that would dilute its emissions impact
% are quantitatively absent in Belgium. Costs were paid forward, but
% customers did not substitute away.

% ----------------------------------------------------------------
% 6.2  WHY BELGIUM DIFFERS FROM CdGM'S FRANCE
% ----------------------------------------------------------------
\subsection{Why Belgium Differs From France}\label{sec:discussion:bevsfr}

% TODO: candidate explanations for the leakage divergence.
%
%   - Tradability of the regulated bundle. Belgian regulated-NACE imports
%     are concentrated in goods (chemicals, basic metals, glass, cement)
%     whose intra-EU sourcing is structurally tight (logistics costs,
%     just-in-time, certification). Substitution toward non-ETS sources
%     would require breaking long-running supplier relationships.
%
%   - Pre-existing intensive intra-EU integration. Belgium's import bundle
%     is dominated by intra-EU flows; the non-ETS share is small at
%     baseline. Even sizeable proportional shifts would deliver small
%     levels.
%
%   - Pass-through abroad. If non-ETS suppliers also pay carbon-equivalent
%     compliance costs (national carbon taxes, RGGI, China ETS pilots),
%     the price wedge between ETS and non-ETS sources is smaller than the
%     headline EUA price suggests.
%
%   - Identification timing. CdGM's window includes the 2005–08 free-allocation
%     era and the 2013+ scarcity era together; Belgian shock is concentrated
%     in Phase IV (2021+) and therefore largely outside CdGM's window.

% ----------------------------------------------------------------
% 6.3  CAVEATS AND OPEN QUESTIONS
% ----------------------------------------------------------------
\subsection{Caveats}\label{sec:discussion:caveats}

% TODO:
%   - Phase IV identification: CPShock fails (macro contamination); CdGM
%     event study with NACE 4d trends works. The cross-section identification
%     of leakage in Phase IV is therefore weaker than in Phases 2–3.
%   - Selection on extensive margin: our pre-2005 leads test for B2B fails
%     parallel trends on the binary spec; the continuous spec has no
%     pre-trend kink but also no Phase-3 effect, so the two diagnostic and
%     null findings are mutually consistent.
%   - Belgium is one country. External validity to e.g. Germany, Italy,
%     Poland is not guaranteed; we discuss what features of the Belgian
%     setting are likely to generalize.
%   - Outcome margins not measured: emissions per unit output (abatement
%     channel) and final-consumer prices (CPI pass-through) are downstream
%     of this paper and addressed elsewhere in the project's pipeline.
